%
% Important note: This latex document should be generated by pdflatex
% to get pdf support right.
%


%%%%%%%%%%%%%%%%%%%%%%%%%%%%%%%%%%%%%%%%%%%%%%%%%%%%%%%%%%%%%%%%%%%%%%%%%%%%%%%%%%%%%%%%%%%%%%%%%%
%%%%%%%%%%% Document class and package options
%%%%%%%%%%%%%%%%%%%%%%%%%%%%%%%%%%%%%%%%%%%%%%%%%%%%%%%%%%%%%%%%%%%%%%%%%%%%%%%%%%%%%%%%%%%%%%%%%%
\documentclass[a4paper,11pt]{report}
\usepackage{graphicx}
\usepackage{hyperref}
\usepackage{fancyvrb}
\usepackage{paralist}
\usepackage{makeidx}
\usepackage{color,calc}
\definecolor{shade}{gray}{0.8}
\newcommand{\makeshadebox}[1]{\fcolorbox{black}{shade}{\makebox{#1}}}


\newcounter{actioncounter}
\addtocounter{actioncounter}{1}

\newcommand{\egads}{\textsc{egads} }
\newcommand{\Egads}{\textsc{Egads} }


\newcommand{\mf}{M\'et\'eoFrance}

\renewcommand{\deg}{$^{\circ}$}

\newcommand{\algdesc}[9]{
\section{#1}\bigskip

\noindent \textbf{Algorithm name:} #2 \bigskip

\noindent \hangindent=2.0cm \textbf{Category:} #4 \bigskip

\noindent \hangindent=2.3cm \textbf{Summary:} #3 \bigskip

\noindent \begin{tabular}{p{2cm} p{2cm} p{8cm}}
	\multicolumn{3}{l}{\textbf{Inputs:}}\\
	\hline
	#5
\end{tabular}

\vspace{0.5cm}


\noindent \begin{tabular}{p{2cm} p{2cm} p{8cm}}
	\multicolumn{3}{l}{\textbf{Outputs:}}\\
	\hline
	#6
\end{tabular}

\vspace{1cm}

\noindent \textbf{Formula:} #7
\bigskip

\noindent \hangindent=1.9cm \textbf{Source:} #8 \bigskip

\noindent \hangindent=2.5cm \textbf{References:}  #9

\pagebreak

}




%%%%%%%%%%%%%%%%%%%%%%%%%%%%%%%%%%%%%%%%%%%%%%%%%%%%%%%%%%%%%%%%%%%%%%%%%%%%%%%%%%%%%%%%%%%%%%%%%%
%%%%%%%%%%% Include page set-up options
%%%%%%%%%%%%%%%%%%%%%%%%%%%%%%%%%%%%%%%%%%%%%%%%%%%%%%%%%%%%%%%%%%%%%%%%%%%%%%%%%%%%%%%%%%%%%%%%%%
%%%%%%%%%%%%%%%%%%%%%%%%%%%%%%%%%%%%%%%%%%%%%%%%%%%%%%%%%%%%%%%%%%%%%%%%%%%%%%%%%%%%%%%%%%%%%%%%%%
%%%%%%%%%%% Page style
%%%%%%%%%%%%%%%%%%%%%%%%%%%%%%%%%%%%%%%%%%%%%%%%%%%%%%%%%%%%%%%%%%%%%%%%%%%%%%%%%%%%%%%%%%%%%%%%%%
\usepackage{fancyhdr}
\setlength{\headheight}{0.4in}
\setlength{\headwidth}{15cm}
\fancyhead[R]{\slshape \leftmark}%  right
%\fancyhead[C]{\textbf{Draft}}
\fancyhead[L]{ % left
	\includegraphics[width=2.5cm]{cover/eufar_a_sm} 
}
\fancyfoot{}
\fancyfoot[R]{\thepage} 
\pagestyle{fancy}
%\pagestyle{headings}


%%%%%%%%%%%%%%%%%%%%%%%%%%%%%%%%%%%%%%%%%%%%%%%%%%%%%%%%%%%%%%%%%%%%%%%%%%%%%%%%%%%%%%%%%%%%%%%%%%
%%%%%%%%%%% Set-up for pdf generation
%%%%%%%%%%%%%%%%%%%%%%%%%%%%%%%%%%%%%%%%%%%%%%%%%%%%%%%%%%%%%%%%%%%%%%%%%%%%%%%%%%%%%%%%%%%%%%%%%%
\hypersetup{
	pdfauthor = {EUFAR},
%	pdftitle = {A \LaTeX Template for Academic Reports},
%	pdfsubject = {Report Template},
	pdfkeywords = {EGADS, algorithms, airborne science},
	pdfcreator = {LaTeX with hyperref package},
	pdfproducer = {latexpdf}
}


%%%%%%%%%%%%%%%%%%%%%%%%%%%%%%%%%%%%%%%%%%%%%%%%%%%%%%%%%%%%%%%%%%%%%%%%%%%%%%%%%%%%%%%%%%%%%%%%%%
%%%%%%%%%%% Page margins
%%%%%%%%%%%%%%%%%%%%%%%%%%%%%%%%%%%%%%%%%%%%%%%%%%%%%%%%%%%%%%%%%%%%%%%%%%%%%%%%%%%%%%%%%%%%%%%%%%
\setlength{\topmargin}{-1cm}
\setlength{\oddsidemargin}{0.5cm}
\setlength{\evensidemargin}{0.5cm}

\setlength{\textwidth}{15cm}
\setlength{\textheight}{22cm}

\setlength{\headsep}{1cm}

%%%%%%%%%%%%%%%%%%%%%%%%%%%%%%%%%%%%%%%%%%%%%%%%%%%%%%%%%%%%%%%%%%%%%%%%%%%%%%%%%%%%%%%%%%%%%%%%%%
%%%%%%%%%%% END OF FILE
%%%%%%%%%%%%%%%%%%%%%%%%%%%%%%%%%%%%%%%%%%%%%%%%%%%%%%%%%%%%%%%%%%%%%%%%%%%%%%%%%%%%%%%%%%%%%%%%%%

\makeindex
\begin{document}
\renewcommand{\thepage}{\roman{page}}

%%%%%%%%%%%%%%%%%%%%%%%%%%%%%%%%%%%%%%%%%%%%%%%%%%%%%%%%%%%%%%%%%%%%%%%%%%%%%%%%%%%%%%%%%%%%%%%%%%
%%%%%%%%%%% Title Page
%%%%%%%%%%%%%%%%%%%%%%%%%%%%%%%%%%%%%%%%%%%%%%%%%%%%%%%%%%%%%%%%%%%%%%%%%%%%%%%%%%%%%%%%%%%%%%%%%%
%%%%%%%%%%%%%%%%%%%%%%%%%%%%%%%%%%%%%%%%%%%%%%%%%%%%%%%%%%%%%%%%%%%%%%%%%%%%%%%%%%%%%%%%%%%%%%%%%%
%%%%%%%%%%% This cover is based on a template provided by LASEC for its semester projects
%%%%%%%%%%%%%%%%%%%%%%%%%%%%%%%%%%%%%%%%%%%%%%%%%%%%%%%%%%%%%%%%%%%%%%%%%%%%%%%%%%%%%%%%%%%%%%%%%%

% Style definitions and re-definitions
\newcommand{\logoeufar}[0]{\noindent\includegraphics[width=2.5cm]{cover/eufar_a_sm}\vspace{0.2cm}\hrule}
\newcommand{\logolabos}[0] {\begin{center}\includegraphics[width=4.5cm]{cover/labos}\end{center}}

\renewcommand{\title}[4]{
\begin{center}
\vspace{4cm}
\textbf{\LARGE{#1}}\\\vspace{.5cm}
\textbf{\large{#2}}\\\vspace{2.5cm}
\textbf{\Huge{#3}}\\\vspace{.5cm}
\textbf{\large{#4}}\end{center}\vspace{.2cm}}

\renewcommand{\author}[2]{\textsf{\begin{center}\Large{#1}\\ \large{#2}\end{center}\vspace{0.3cm}}}
\newcommand{\meetingdate}[2]{\begin{center}\large{#1.  #2.}\end{center}}

\renewcommand{\date}[2]{\begin{center}\normalsize{Last Updated: #1}\end{center}}


\thispagestyle{empty}

\setlength{\voffset}{.0cm}

\logoeufar
\title{EUFAR FP7}{N6SP - Standards and Protocols}{EGADS Algorithm Handbook}{Version 0.5.6}
%\author{}{}


\vspace{4.0cm}
\textsf{\date{\today}}


\vspace{4cm}
\hrule
%%%%%%%%%%%%%%%%%%%%%%%%%%%%%%%%%%%%%%%%%%%%%%%%%%%%%%%%%%%%%%%%%%%%%%%%%%%%%%%%%%%%%%%%%%%%%%%%%%
%%%%%%%%%%% END OF FILE
%%%%%%%%%%%%%%%%%%%%%%%%%%%%%%%%%%%%%%%%%%%%%%%%%%%%%%%%%%%%%%%%%%%%%%%%%%%%%%%%%%%%%%%%%%%%%%%%%%





%%%%%%%%%%%%%%%%%%%%%%%%%%%%%%%%%%%%%%%%%%%%%%%%%%%%%%%%%%%%%%%%%%%%%%%%%%%%%%%%%%%%%%%%%%%%%%%%%%
%%%%%%%%%%% Indices
%%%%%%%%%%%%%%%%%%%%%%%%%%%%%%%%%%%%%%%%%%%%%%%%%%%%%%%%%%%%%%%%%%%%%%%%%%%%%%%%%%%%%%%%%%%%%%%%%%
\tableofcontents
\pagebreak


%%%%%%%%%%%%%%%%%%%%%%%%%%%%%%%%%%%%%%%%%%%%%%%%%%%%%%%%%%%%%%%%%%%%%%%%%%%%%%%%%%%%%%%%%%%%%%%%%%
%%%%%%%%%%% Introduction
%%%%%%%%%%%%%%%%%%%%%%%%%%%%%%%%%%%%%%%%%%%%%%%%%%%%%%%%%%%%%%%%%%%%%%%%%%%%%%%%%%%%%%%%%%%%%%%%%%
\renewcommand{\thepage}{\arabic{page}}
\setcounter{page}{1}

%%%%%%%%%%%%%%%%%%%%%%%%%%%%%%%%%%%%%%%%%%%%%%%%%%%%%%%%%%%%%%%%%%%%%%%%%%%%%%%%%%%%%%%%%%%%%%%%%%
%%%%%%%%%%% Chapters
%%%%%%%%%%%%%%%%%%%%%%%%%%%%%%%%%%%%%%%%%%%%%%%%%%%%%%%%%%%%%%%%%%%%%%%%%%%%%%%%%%%%%%%%%%%%%%%%%%

\chapter{Introduction}

This document contains descriptions of algorithms contained in the \egads toolbox. Within each algorithm description is the following:

\begin{itemize}
\item \textbf{Algorithm Name} -- name of algorithm as implemented in \egads.
\item \textbf{Category} -- general category of algorithm. Algorithm can be found in this subdirectory in \egads.
\item \textbf{Summary} -- short description of what the algorithm does.
\item \textbf{Inputs} -- expected inputs to algorithm. This field includes expected units, and data type of input.
\item \textbf{Outputs} -- outputs produced by algorithm.
\item \textbf{Formula} -- description of formulas or methods behind the algorithm.
\item \textbf{Source} -- person, institution or entity who provided the algorithm.
\item \textbf{References} -- any references to literature, journals or documents with more information on the current algorithm
\end{itemize}

To aid in algorithm usage, there is a general naming scheme for \egads algorithms. Generally, algorithm names are composed as follows: |{general calculation}_{specific calculation}_{source}|. So, for example, an algorithm provided by CNRM to calculate the density of dry air would be named |density_dry_air_cnrm|.


%\chapter{Spatial Transformations}

%\chapter{Temporal Transformations}

\chapter{Thermodynamics}

%% $Date$
%% $Revision$

\algdesc{Pressure altitude
}
{ %%%%%% Algorithm name %%%%%%
altitude\_pressure\_cnrm
}
{ %%%%%% Algorithm summary %%%%%%
Calculates pressure altitude using virtual temperature.
}
{ %%%%%% Category %%%%%%
Thermodynamics
}
{ %%%%%% Inputs %%%%%%
$T_v$ & Vector & Virtual temperature [K or \deg C] \\
$P_s$ & Vector & Static pressure [hPa] \\
$P_{surface}$ & Coeff & Surface pressure [hPa] \\ 
$R_a/g$ & Coeff & Gas constant of air over acceleration of gravity
}
{ %%%%%% Outputs %%%%%%
$Alt_p$ & Vector & Pressure altitude [m]
}
{ %%%%%% Formula %%%%%%
\begin{displaymath}
 Alt_p = \frac{R_a}{g} T_v \log\left(\frac{P_{surface}}{P_s}\right)
\end{displaymath}

}
{ %%%%%% Author %%%%%%
CNRM/GMEI/TRAMM
}
{ %%%%%% References %%%%%% 

}




%% $Date$
%% $Revision$
\index{density\_dry\_air\_cnrm}
\algdesc{Density of dry air}
{ %%%%%% Algorithm name %%%%%%
density\_dry\_air\_cnrm
}
{ %%%%%% Algorithm summary %%%%%%
Calculates density of dry air given static temperature and pressure.
}
{ %%%%%% Category %%%%%%
Thermodynamics
}
{ %%%%%% Inputs %%%%%%
$P_s$ & Vector & Static pressure [hPa] \\
$T_s$ & Vector & Static temperature [K or \deg C]
}
{ %%%%%% Outputs %%%%%%
$\rho$ & Vector & Density of dry air [kg/m$^3$]
}
{ %%%%%% Formula %%%%%%
\begin{displaymath}
 \rho = \frac{100 P_s}{R_a T_s}
\end{displaymath}
%
with $R_a = 287.05$ J kg$^-1$ K$^-1$ 

Density of humid air can be calculated using this same algorithm by using virtual temperature instead of static temperature.
}
{ %%%%%% Author %%%%%%
CNRM/GMEI/TRAMM
}
{ %%%%%% References %%%%%% 
Equation d'\'etat d'un gaz parfait, Triplet-Roche \cite{Triplet}, page 34.
}




%% $Date$
%% $Revision$

\algdesc{Relative humidity from capacitive probe
}
{ %%%%%% Algorithm name %%%%%%
hum\_rel\_capacitive\_cnrm
}
{ %%%%%% Algorithm summary %%%%%%
Calculates relative humidity using the measured frequency from a capacitive probe.
}
{ %%%%%% Inputs %%%%%%
$Ucapf$ & Vector & Output frequency of the capacitive probe [Hz]\\ 
$T_s$ & Vector & Static temperature [K] \\
$P_s$ & Vector & Static pressure [hPa] \\
$\Delta P$ & Vector & Dynamic pressure [hPa] \\
$C_t$ & Coeff. & Temperature correction coefficient [\%$^\circ$C] \\
$F_{min}$ & Coeff. & Minimal acceptable frequency [Hz] \\
$C_0$ & Coeff. & 0th degree calibration coefficient \\
$C_1$ & Coeff. & 1st degree calibration coefficient \\
$C_2$ & Coeff. & 2nd degree calibration coefficient 
}
{ %%%%%% Outputs %%%%%%
$H_u$ & Vector & Relative humidity [\%]
}
{ %%%%%% Formula %%%%%%
If $Ucapf \leq F_{min}$ then $Ucapf = F_{min}$

\begin{displaymath}
H_u = \frac{P_s}{P_s + \Delta P} \left[C_0 + C_1 Ucapf + C_2 Ucapf^2 + C_t (T_s-20)\right] \nonumber
\end{displaymath}
with $T_s$ in $^\circ C$.
}
{ %%%%%% Author %%%%%%
CNRM/GMEI/TRAMM
}
{ %%%%%% References %%%%%% 
H.~Bellec and G.~Duverneuil.  Appareils de mesure de l'hygrom\'etrie sur le Merlin IV.  Note de Centre 9, M\'et\'eo-France CNRM/CAM, July 1996.
}

%% $Date$
%% $Revision$

\algdesc{Pressure and angle of incidence}
{ %%%%%% Algorithm name %%%%%%
pressure\_angle\_incidence\_cnrm
}
{ %%%%%% Algorithm summary %%%%%%
Calculates static pressure and dynamic pressure by correction of static error. Angle of attack and sideslip are calculated from the horizontal and vertical differential pressures.
}
{ %%%%%% Category %%%%%%
Thermodynamics
}
{ %%%%%% Inputs %%%%%%
$P_{sr}$ & Vector & Raw static pressure [hPa] \\
$\Delta P_r$ & Vector & Raw dynamic pressure [hPa] \\
$\Delta P_h$ & Vector & Horizontal differential pressure [hPa] \\
$\Delta P_v$ & Vector & Vertical differential pressure [hPa] \\
$C_\alpha$ & Coeff.[2] & Angle of attack calibration coefficients \\
$C_\beta$ & Coeff.[2] & Slip calibration coefficients \\
$C_{errstat}$ & Coeff.[4] & Static error coefficients 
}
{ %%%%%% Outputs %%%%%%
$P_s$ & Vector & Static Pressure [hPa] \\
$\Delta P$ & Vector & Dynamic pressure corrected with static error [hPa] \\
$\alpha$ & Vector & Angle of attack [rad] \\
$\beta$ & Vector & Sideslip [rad] 
}
{ %%%%%% Formula %%%%%%
If $\Delta P_r > 25$hPa:
\begin{displaymath}
 Errstat = C_{errstat}[0] + C_{errstat}[1] \Delta P_r + C_{errstat}[2] \Delta P_r^2 + C_{errstat}[3] \Delta P_r^3
\end{displaymath}
%
otherwise:
\begin{eqnarray}
 Errstat = \frac{\Delta P_r}{25} \textrm{ Errstat @ 25 hPa} \\ \nonumber
 P_s = P_{sr} - Errstat \\ \nonumber
 \Delta P = \Delta P_r + Errstat \\ \nonumber
 \alpha = C_\alpha [0] + C_\alpha [1] \frac{\Delta P_v}{\Delta P} \\ \nonumber
 \beta = C_\beta [0] + C_\beta [1] \frac{\Delta P_h}{\Delta P} \nonumber
\end{eqnarray}
}
{ %%%%%% Author %%%%%%
CNRM/GMEI/TRAMM
}
{ %%%%%% References %%%%%% 

}




%% $Date$
%% $Revision$
\index{temp\_potential\_cnrm}
\algdesc{Potential Temperature}
{ %%%%%% Algorithm name %%%%%%
temp\_potential\_cnrm
}
{ %%%%%% Algorithm summary %%%%%%
Calculates potential temperature.
}
{ %%%%%% Category %%%%%%
Thermodynamics
}
{ %%%%%% Inputs %%%%%%
$T_s$ & Vector & Static temperature [K or $\circ$C] \\
$P_s$ & Vector & Static pressure [hPa] \\
$R_a/c_{pa}$ & Coeff. & Gas constant of air divided by specific heat of air at constant pressure
}
{ %%%%%% Outputs %%%%%%
$\theta$ & Vector & Potential temperature [same unit as $T_s$]
}
{ %%%%%% Formula %%%%%%
\begin{displaymath}
 \theta = T_s \left(\frac{1000}{P_s}\right)^{R_a/c_{pa}}
\end{displaymath}
}
{ %%%%%% Author %%%%%%
CNRM/GMEI/TRAMM
}
{ %%%%%% References %%%%%% 
Triplet-Roche \cite{Triplet}.
}




%% $Date$
%% $Revision$

\algdesc{Equivalent Potential Temperature}
{ %%%%%% Algorithm name %%%%%%
temp\_potential\_equiv\_cnrm
}
{ %%%%%% Algorithm summary %%%%%%
Calculates equivalent potential temperature of air. The equivalent potential temperature is the temperature a parcel of air would reach
if all water vapor in the parcel condensed, and the parcel was brought adiabatially to 1000 hPa.
}
{ %%%%%% Category %%%%%%
Thermodynamics
}
{ %%%%%% Inputs %%%%%%
$T_s$ & Vector & Static temperature [K or $\circ$C] \\
$\theta$ & Vector & Potential temperature [K or $\circ$C] \\
$r$ & Vector & Vater vapor mixing ratio [g/kg] \\
$c_{pa}$ & Coeff. & Specific heat of dry air at constant pressure
}
{ %%%%%% Outputs %%%%%%
$\theta_e$ & Vector & Equivalent potential temperature [same units as $T_s$]
}
{ %%%%%% Formula %%%%%%
\begin{displaymath}
 \theta_e = \theta \left(1 + r \frac{L}{c_{pa} T_s} \right)
\end{displaymath}
%
where $L = 3136.17 - 2.34 T_s$ (for $T_s$ in K)
}
{ %%%%%% Author %%%%%%
CNRM/GMEI/TRAMM
}
{ %%%%%% References %%%%%% 
From the CAM routine which is identical to the algorithm P.~Durand cited in the formula book created for PYREX.
}




%% $Date$
%% $Revision$

\algdesc{Static Temperature}
{ %%%%%% Algorithm name %%%%%%
temp\_static\_cnrm \index{temp_static_cnrm}
}
{ %%%%%% Algorithm summary %%%%%%
Calculates static temperature of the air from total temperature.  This method applies to probe types such as the Rosemount.
}
{ %%%%%% Category %%%%%%
Thermodynamics
}
{ %%%%%% Inputs %%%%%%
$T_t$ &	Vector & Measured total temperature [K] \\
${\Delta}P$ & Vector & Dynamic pressure [hPa] \\
$P_s$ & Vector & Static pressure [hPa] \\
$r_f$ & Coeff. & Probe recovery coefficient \\ 
$R_a/c_{pa}$ & Coeff. & Gas constant of air divided by specific heat of air at constant pressure
}
{ %%%%%% Outputs %%%%%%
$T_s$ & Vector & Static temperature [K]
}
{ %%%%%% Formula %%%%%%
\begin{displaymath}
 T_s = \frac{T_t}{1+r_f \left(\left(1+\frac{\Delta P}{P_s}\right)^{R_a/c_{pa}}-1\right)} \nonumber
\end{displaymath}
}
{ %%%%%% Author %%%%%%
CNRM/GMEI/TRAMM
}
{ %%%%%% References %%%%%%
}


%% $Date$
%% $Revision$
\index{temp\_virtual\_cnrm}
\algdesc{Virtual Temperature}
{ %%%%%% Algorithm name %%%%%%
temp\_virtual\_cnrm
}
{ %%%%%% Algorithm summary %%%%%%
Calculates the virtual temperature of air.
}
{ %%%%%% Category %%%%%%
Thermodynamics
}
{ %%%%%% Inputs %%%%%%
$T_s$ & Vector & Static temperature [K or $\circ$C] \\
$r$ & Vector & Water vapor mixing ratio [g/kg]
}
{ %%%%%% Outputs %%%%%%
$T_v$ & Vector & Virtual temperature [same units as $T_s$]
}
{ %%%%%% Formula %%%%%%
\begin{displaymath}
 T_v = T_s \frac{1 + (R_v/R_a) r }{1 + r}
\end{displaymath}
%
where $R_v/R_a = 1.608$ 
}
{ %%%%%% Author %%%%%%
CNRM/GMEI/TRAMM
}
{ %%%%%% References %%%%%% 
Triplet-Roche \cite{Triplet}, page 56.
}




%% $Date: 2012-07-06 17:42:54#$
%% $Revision$
\index{velocity\_tas\_cnrm}
\algdesc{True air speed}
{ %%%%%% Algorithm name %%%%%%
velocity\_tas\_cnrm
}
{ %%%%%% Algorithm summary %%%%%%
Calculates true air speed based on static pressure, static temperature and dynamic pressure using the Barr\'e-St Venant formula.  
}
{ %%%%%% Category %%%%%%
Thermodynamics
}
{ %%%%%% Inputs %%%%%%
$T_s$ &	Vector & Static temperature [K] \\
${\Delta}P$ & Vector & Dynamic pressure [hPa] \\
$P_s$ & Vector & Static pressure [hPa] \\
$c_{pa}$ & Coeff. & Specific heat of air at constant pressure (for dry air 1004 J K$^{-1}$ kg$^{-1}$) \\ 
$R_a/c_{pa}$ & Coeff. & Gas constant of air divided by specific heat of air at constant pressure
}
{ %%%%%% Outputs %%%%%%
$V_t$ & Vector & True air speed [m/s]
}
{ %%%%%% Formula %%%%%%
\begin{displaymath}
V_t^2 = 2 c_{pa} T_s \left[\left(1+\frac{\Delta P}{P_s}\right)^{R_a/c_{pa}}-1\right]
\end{displaymath}
}
{ %%%%%% Author %%%%%%
CNRM/GMEI/TRAMM
}
{ %%%%%% References %%%%%% 
NCAR-RAF Bulletin \#23 \cite{NCAR23}, \textit{M\'echanique des fluides}, Candel \cite{Candel} 
}




%% $Date$
%% $Revision$

\algdesc{Longitudinal true airspeed}
{ %%%%%% Algorithm name %%%%%%
velocity\_tas\_longitudinal\_cnrm
}
{ %%%%%% Algorithm summary %%%%%%
Calculates the true air speed along the longitudinal axis of the aircraft.
}
{ %%%%%% Inputs %%%%%%
$V_t$ & Vector & True air speed [m/s] \\
$\alpha$ & Vector & Angle of attack [rad] \\
$\beta$ & Vector & Sideslip angle [rad] 
}
{ %%%%%% Outputs %%%%%%
$V_{tx}$ & Vector & Longitudinal true air speed [m/s]
}
{ %%%%%% Formula %%%%%%
\begin{displaymath}
 V_{tx} = \frac{V_t}{\sqrt{1 + \tan^2 \alpha + \tan^2 \beta}}
\end{displaymath}
}
{ %%%%%% Author %%%%%%
CNRM/GMEI/TRAMM
}
{ %%%%%% References %%%%%% 
NCAR-RAF Bulletin \#23 \cite{NCAR25}
}






\chapter{Microphysics}

%% $Date$
%% $Revision$
\index{mean\_diameter\_raf}
\algdesc{Mean diameter}
{ %%%%%% Algorithm name %%%%%%
mean\_diameter\_raf
}
{ %%%%%% Algorithm summary %%%%%%
Calculates the arithmetic average of all particle diameters given in a particle size distribution.
}
{ %%%%%% Category %%%%%%
Microphysics
}
{ %%%%%% Inputs %%%%%%
$n_i$ & Array[time, bins] & Number of particles in each channel $i$ \\
$d_i$ & Vector[bins] & Channel $i$ size [$\mu$m]
}
{ %%%%%% Outputs %%%%%%
$\bar{D}$ & Vector[time] & Mean diameter [$\mu$m]
}
{ %%%%%% Formula %%%%%%
\begin{displaymath}
 \bar{D} = \frac{\sum_i n_i d_i}{N_t}
\end{displaymath}
where $N_t$ is the total number of particles.
}
{ %%%%%% Author %%%%%%
NCAR-RAF
}
{ %%%%%% References %%%%%% 
NCAR-RAF Bulletin No. 24 -- \href{http://www.eol.ucar.edu/raf/Bulletins/bulletin24.html}{http://www.eol.ucar.edu/raf/Bulletins/bulletin24.html}
}



%% $Date: 2012-07-06 17:42:54#$
%% $Revision$
\index{number\_conc\_total\_raf}
\algdesc{Total Number Concentration}
{ %%%%%% Algorithm name %%%%%%
number\_conc\_total\_raf
}
{ %%%%%% Algorithm summary %%%%%%
Calculation of total number concentration for a particle probe.
}
{ %%%%%% Category %%%%%%
Microphysics
}
{ %%%%%% Inputs %%%%%%
$n_i$ & Array & Number of particles in each channel $i$ \\
$SV$ & Array & Sample volume [m$^3$] 
}
{ %%%%%% Outputs %%%%%%
$N_t$ & Vector & Total number concentration [m$^{-3}$]
}
{ %%%%%% Formula %%%%%%
\begin{displaymath}
 N_t = \frac{\sum_{i} n_i}{SV}
\end{displaymath}

}
{ %%%%%% Author %%%%%%
NCAR-RAF
}
{ %%%%%% References %%%%%% 
NCAR-RAF Bulletin No. 24. \cite{NCAR24}
}



%% $Date: 2012-07-06 17:42:54#$
%% $Revision$
\index{sample\_area\_oap\_all\_in\_raf}
\algdesc{Sample area for imaging probes}
{ %%%%%% Algorithm name %%%%%%
sample\_area\_oap\_all\_in\_raf
}
{ %%%%%% Algorithm summary %%%%%%
Calculation of 'all in' sample area size for OAP probes such as the 2DC, 2DP, CIP, etc. This sample area varies by
number of shadowed diodes. This routine calculates a sample area per bin.
}
{ %%%%%% Category %%%%%%
Microphysics
}
{ %%%%%% Inputs %%%%%%
$\lambda$ & Coeff. & Laser wavelength [nm] \\
$D_{arms}$ & Coeff. & Distance between probe arm tips [mm] \\
dD & Coeff. & Diode diameter [$\mu$m] \\
M & Coeff. & Probe magnification factor \\
N & Coeff. & Number of diodes in array
}
{ %%%%%% Outputs %%%%%%
SA & Vector & Sample area [m$^2$]
}
{ %%%%%% Formula %%%%%%
\begin{eqnarray}
DOF_i = \frac{6 R_i^2}{\lambda} \\ \nonumber
R_i = i \frac{dD}{2} \\ \nonumber
X = {1...N} \nonumber
\end{eqnarray}
where $DOF$ must be less than $D_{arms}$. The parameter $i$ ranges from 1 to $N$. 


\begin{displaymath}
ESW_i = \frac{dD(N-X_i-1)}{M}
\end{displaymath}
A value for $ESW_i$ (effective sample width) is calculated for each $X$.

\begin{displaymath}
 SA_i = (DOF_i)(ESW_i) 
\end{displaymath}
}
{ %%%%%% Author %%%%%%
NCAR-RAF
}
{ %%%%%% References %%%%%% 
NCAR-RAF Bulletin No. 24. \cite{NCAR24}
}



%% $Date$
%% $Revision$
\index{sample\_area\_oap\_center\_in\_raf}
\algdesc{Sample area for imaging probes}
{ %%%%%% Algorithm name %%%%%%
sample\_area\_oap\_center\_in\_raf
}
{ %%%%%% Algorithm summary %%%%%%
Calculation of 'center in' sample area size for OAP probes such as the 2DC, 2DP, CIP, etc. This sample area varies by number of shadowed diodes. This routine is intended to calculate a sample area per bin.
}
{ %%%%%% Category %%%%%%
Microphysics
}
{ %%%%%% Inputs %%%%%%
$\lambda$ & Coeff. & Laser wavelength [nm] \\
$D_{arms}$ & Coeff. & Distance between probe arm tips [mm] \\
dD & Coeff. & Diode diameter [$\mu$m] \\
M & Coeff. & Probe magnification factor \\
N & Coeff. & Number of diodes in array \\
}
{ %%%%%% Outputs %%%%%%
SA & Vector & Sample area [m$^2$]
}
{ %%%%%% Formula %%%%%%
\begin{eqnarray}
DOF_i = \frac{6 R_i^2}{\lambda} \\ \nonumber
R_i = X \frac{dD}{2} \\ \nonumber
X = {1...N} \nonumber
\end{eqnarray}
where $DOF$ must be less than $D_{arms}$. The parameter $i$ ranges from 1 to $N$.

\begin{displaymath}
ESW = N dD \nonumber
\end{displaymath}

\begin{displaymath}
 SA_i = (DOF_i)(ESW) 
\end{displaymath}
}
{ %%%%%% Author %%%%%%
NCAR-RAF
}
{ %%%%%% References %%%%%% 
NCAR-RAF Bulletin No. 24. \cite{NCAR24}
}



%% $Date: 2010-09-17 17:53:24 +0200 (Fri, 17 Sep 2010) $
%% $Revision: 24 $
\index{sample\_area\_scattering\_raf}
\algdesc{Sample area for scattering probes}
{ %%%%%% Algorithm name %%%%%%
sample\_area\_scattering\_raf
}
{ %%%%%% Algorithm summary %%%%%%
Calculation of sample area for scattering probes such as the FSSP, CAS, etc.
}
{ %%%%%% Category %%%%%%
Microphysics
}
{ %%%%%% Inputs %%%%%%
DOF & Coeff. & Depth of field [m] \\
BD & Coeff. & Beam diameter [m]
}
{ %%%%%% Outputs %%%%%%
SA & Coeff. & Sample area [m$^2$]
}
{ %%%%%% Formula %%%%%%
\begin{displaymath}
SA = (DOF) (BD)
\end{displaymath}
}
{ %%%%%% Author %%%%%%
NCAR-RAF
}
{ %%%%%% References %%%%%% 
NCAR-RAF Bulletin No. 24 -- \href{http://www.eol.ucar.edu/raf/Bulletins/bulletin24.html}{http://www.eol.ucar.edu/raf/Bulletins/bulletin24.html}
}



%% $Date$
%% $Revision$
\index{sample\_volume\_general\_raf}
\algdesc{Sample Volume}
{ %%%%%% Algorithm name %%%%%%
sample\_volume\_general\_raf
}
{ %%%%%% Algorithm summary %%%%%%
Calculates sample volume for microphysics probes (1D, 2D, FSSP, etc).
}
{ %%%%%% Category %%%%%%
Microphysics
}
{ %%%%%% Inputs %%%%%%
$V_t$ & Vector & True air speed [m/s] \\
SA & Coeff. & Sample area of probe [m$^2$] \\
$t_s$ & Coeff. & Sample rate [s] \\
}
{ %%%%%% Outputs %%%%%%
SV & Vector & Sample volume [m$^3$]
}
{ %%%%%% Formula %%%%%%
\begin{displaymath}
 SV = V_t t_s SA
\end{displaymath}
}
{ %%%%%% Author %%%%%%
NCAR-RAF
}
{ %%%%%% References %%%%%% 
NCAR-RAF Bulletin No. 24 -- \href{http://www.eol.ucar.edu/raf/Bulletins/bulletin24.html}{http://www.eol.ucar.edu/raf/Bulletins/bulletin24.html}
}





\printindex

%%%%%%%%%%%%%%%%%%%%%%%%%%%%%%%%%%%%%%%%%%%%%%%%%%%%%%%%%%%%%%%%%%%%%%%%%%%%%%%%%%%%%%%%%%%%%%%%%%
%%%%%%%%%%% APPENDIX inclusions
%%%%%%%%%%%%%%%%%%%%%%%%%%%%%%%%%%%%%%%%%%%%%%%%%%%%%%%%%%%%%%%%%%%%%%%%%%%%%%%%%%%%%%%%%%%%%%%%%%
%\begin{appendix}
%\include{02_appendices/01/appendix_01}
%\end{appendix}



%%%%%%%%%%%%%%%%%%%%%%%%%%%%%%%%%%%%%%%%%%%%%%%%%%%%%%%%%%%%%%%%%%%%%%%%%%%%%%%%%%%%%%%%%%%%%%%%%%
%%%%%%%%%%% Bibliography
%%%%%%%%%%%%%%%%%%%%%%%%%%%%%%%%%%%%%%%%%%%%%%%%%%%%%%%%%%%%%%%%%%%%%%%%%%%%%%%%%%%%%%%%%%%%%%%%%%
%\newpage\addcontentsline{toc}{section}{Bibliography and references}
%\begin{thebibliography}{5}
%\bibitem{Vahdat et al., 2002} Vahdat A., Yocum K., Walsh K., Mahadevan P., Kosti\c{c} D., Chase J., and Becker D. (2002). Scalability and Accuracy in a Large-Scale Network Emulator. \emph{Proceedings of 5th Symposium on Operating Systems Design and Implementation (OSDI)}
%\bibitem{Modelnet} Modelnet. http://issg.cs.duke.edu/modelnet.html.
%\bibitem{ModelNetRelease} Modelnet release page. http://sysnet.ecsd.edu/modelnet/realease.html.
%\bibitem{ModelNetHowto} Modelnet Howto. http://sysnet.ucsd.edu/modelnet/howto.html.
%\bibitem{Holenstein, 2005} Holenstein R. (2005). Scalability and Accuracy in a Large-Scale Network Emulator - CPSC 538A - Paper Presentation
%\bibitem{Harpin, 2001} Harpin, T. (2001). Using java.lang.reflect.Proxy to Interpose on Java Class Methods. \emph{Sun Developer Network (SDN)}. http://java.sun.com/developer/technicalArticles/JavaLP/Interposing/.
%\bibitem{Kostic, 2006} Kostic D. (2005). Modelnet notes. http://resolute.ucsd.edu/dokuwiki/notes:modelnet.
%\end{thebibliography}


\end{document}
%%%%%%%%%%%%%%%%%%%%%%%%%%%%%%%%%%%%%%%%%%%%%%%%%%%%%%%%%%%%%%%%%%%%%%%%%%%%%%%%%%%%%%%%%%%%%%%%%%
%%%%%%%%%%% END OF FILE
%%%%%%%%%%%%%%%%%%%%%%%%%%%%%%%%%%%%%%%%%%%%%%%%%%%%%%%%%%%%%%%%%%%%%%%%%%%%%%%%%%%%%%%%%%%%%%%%%%
